\documentclass[../../main.tex]{subfiles}
\begin{document}

\section*{问题}\label{section:ch01-问题}

\begin{exercise}\label{exercise:ch01-1-1}
设 $X$ 为所有满足 $y=\pm 1$ 的点 $(x,y)\in\mathbb{R}^{2}$ 所成集合，令 $M$ 为 $X$ 关于由 $(x,-1)\sim(x,1)$（对所有 $x\neq 0$）生成的等价关系的商空间。证明 $M$ 是局部欧几里得的且第二可数，但不是 Hausdorff 的。（这个空间称为带两原点的直线。）
\end{exercise}

\begin{exercise}\label{exercise:ch01-1-2}
证明：不可数个 $\mathbb{R}$ 的拷贝的不交并是局部欧几里得的且 Hausdorff，但不是第二可数的。
\end{exercise}

\begin{exercise}\label{exercise:ch01-1-3}
称一个拓扑空间是 $\sigma$-紧的，如果它可以表示为可数多个紧子空间的并。证明：一个局部欧几里得的 Hausdorff 空间是拓扑流形，当且仅当它是 $\sigma$-紧的。
\end{exercise}

\begin{exercise}\label{exercise:ch01-1-4}
设 $M$ 是一个拓扑流形，$\mathcal{U}$ 是 $M$ 的一个开覆盖。
\begin{enumerate}[(a)]
\item 假设 $\mathcal{U}$ 中每个集合只与有限多个其他集合相交，证明 $\mathcal{U}$ 是局部有限的。
\item 举一个例子说明 (a) 的逆可能不成立。
\item 现在假设 $\mathcal{U}$ 中的集合在 $M$ 中都是预紧的；证明逆命题：若 $\mathcal{U}$ 是局部有限的，则 $\mathcal{U}$ 中每个集合只与有限多个其他集合相交。
\end{enumerate}
\end{exercise}

\begin{exercise}\label{exercise:ch01-1-5}
假设 $M$ 是一个局部欧几里得的 Hausdorff 空间。证明 $M$ 是第二可数的，当且仅当它是仿紧的且有可数多个连通分支。[提示：假设 $M$ 是仿紧的，证明 $M$ 的每个分支都有一个由预紧坐标域组成的局部有限覆盖，并从中抽取一个可数子覆盖。]
\end{exercise}

\begin{exercise}\label{exercise:ch01-1-6}
设 $M$ 是一个维数 $n\geqslant 1$ 的非空拓扑流形。若 $M$ 有一个光滑结构，证明它有不可数个不同的光滑结构。[提示：首先证明对任意 $s>0$，映射 $F_{s}(x)=|x|^{s-1}x$ 定义一个从 $B^{n}$ 到自身的同胚，且它是微分同胚当且仅当 $s=1$。]
\end{exercise}

\begin{exercise}\label{exercise:ch01-1-7}
令 $N=(0,\dots,0,1)\in S^{n}\subset\mathbb{R}^{n+1}$ 表示北极，$S=(0,\dots,0,-1)$ 表示南极。定义球极投影 $\beta:S^{n}\setminus\{N\}\to\mathbb{R}^{n}$ 为
\begin{align*}
\beta\bigl(x^{1},\dots,x^{n+1}\bigr)=\frac{(x^{1},\dots,x^{n})}{1-x^{n+1}}.
\end{align*}
令 $\widetilde{\beta}(x)=-\beta(-x)$，其中 $x\in S^{n}\setminus\{S\}$。
\begin{enumerate}[(a)]
\item 对任意 $x\in S^{n}\setminus\{N\}$，证明 $\beta(x)=u$，其中 $(u,0)$ 是过 $N$ 与 $x$ 的直线与线性子空间 $x^{n+1}=0$ 的交点（图~\cref{figure:ch01-1.13}）。类似地，证明 $\widetilde{\beta}(x)$ 是过 $S$ 与 $x$ 的直线与同一子空间的交点。（由于这个原因，$\widetilde{\beta}$ 称为从南极出发的球极投影。）
\item 证明 $\beta$ 是双射，且
\begin{align*}
\beta^{-1}\bigl(u^{1},\dots,u^{n}\bigr)=\frac{(2u^{1},\dots,2u^{n},|u|^{2}-1)}{|u|^{2}+1}.
\end{align*}
\item 计算转移映射 $\widetilde{\beta}\circ\beta^{-1}$，并验证由两个图 $(S^{n}\setminus\{N\},\beta)$ 和 $(S^{n}\setminus\{S\},\widetilde{\beta})$ 组成的图册在 $S^{n}$ 上定义了一个光滑结构。（由 $\beta$ 或 $\widetilde{\beta}$ 定义的坐标称为球极坐标。）
\item 证明这个光滑结构与例~\cref{example:ch01-球面光滑结构} 中定义的那个相同。
\end{enumerate}
（用于第~201、269、301、345、347、450 页。）
\end{exercise}

\begin{figure}[ht]
\centering
\includegraphics[width=0.8\textwidth]{fig-01-13.png}
\caption{球极投影}\label{figure:ch01-1.13}
\end{figure}

\begin{exercise}\label{exercise:ch01-1-8}
通过把 $\mathbb{R}^{2}$ 等同于 $\mathbb{C}$，我们可以把单位圆 $S^{1}$ 看作复平面的一个子集。$S^{1}$ 的一个子集 $U$ 上的角函数是指一个连续函数 $\theta:U\to\mathbb{R}$，使对所有 $z\in U$ 有 $e^{\mathrm{i}\theta(z)}=z$。证明：存在 $S^{1}$ 的某个开子集 $U$ 上的角函数 $\theta$，当且仅当 $U\neq S^{1}$。对任意这样的角函数，证明 $(U,\theta)$ 是 $S^{1}$ 带其标准光滑结构的一个光滑坐标图。（用于第~37、152、176 页。）
\end{exercise}

\begin{exercise}\label{exercise:ch01-1-9}
复射影 $n$ 空间，记作 $\mathbb{CP}^{n}$，是 $\mathbb{C}^{n+1}$ 的所有 $1$ 维复线性子空间所成的集合，带有从自然投影 $\pi:\mathbb{C}^{n+1}\setminus\{0\}\to\mathbb{CP}^{n}$ 继承的商拓扑。证明 $\mathbb{CP}^{n}$ 是一个紧的 $2n$ 维拓扑流形，并说明如何赋予它一个与我们为 $\mathbb{RP}^{n}$ 构造的那个类似的光滑结构。（我们利用对应
\begin{align*}
\bigl(x^{1}+\mathrm{i}y^{1},\dots,x^{n+1}+\mathrm{i}y^{n+1}\bigr)\leftrightarrow\bigl(x^{1},y^{1},\dots,x^{n+1},y^{n+1}\bigr)
\end{align*}
把 $\mathbb{C}^{n+1}$ 等同于 $\mathbb{R}^{2n+2}$。）（用于第~48、96、172、560、561 页。）
\end{exercise}

\begin{exercise}\label{exercise:ch01-1-10}
设 $k$ 和 $n$ 是满足 $0<k<n$ 的整数，令 $P,Q\subset\mathbb{R}^{n}$ 分别为由 $(e_{1},\dots,e_{k})$ 和 $(e_{k+1},\dots,e_{n})$ 张成的线性子空间，其中 $e_{i}$ 是 $\mathbb{R}^{n}$ 的第 $i$ 个标准基向量。对任意与 $Q$ 平凡相交的 $k$ 维子空间 $S\subset\mathbb{R}^{n}$，证明例~\cref{example:ch01-Grassmann流形} 中构造的坐标表示 $\varphi(S)$ 是唯一的一个 $(n-k)\times k$ 矩阵 $B$，使 $S$ 由矩阵
\begin{align*}
\begin{pmatrix}I_{k}\\ B\end{pmatrix}
\end{align*}
的列张成，其中 $I_{k}$ 表示 $k\times k$ 单位矩阵。
\end{exercise}

\begin{exercise}\label{exercise:ch01-1-11}
令 $M=\overline{B^{n}}$，即 $\mathbb{R}^{n}$ 中的闭单位球。证明 $M$ 是一个拓扑带边流形，其中 $S^{n-1}$ 中每个点都是边界点，$B^{n}$ 中每个点都是内点。说明如何赋予它一个光滑结构，使每个光滑内图都是 $B^{n}$ 上标准光滑结构的一个光滑图。[提示：考虑映射 $\pi\circ\beta^{-1}:\mathbb{R}^{n}\to\mathbb{R}^{n}$，其中 $\beta:S^{n}\to\mathbb{R}^{n}$ 是球极投影（习题~\cref{exercise:ch01-1-7}），$\pi$ 是从 $\mathbb{R}^{n+1}$ 到 $\mathbb{R}^{n}$ 略去除最后一个坐标外某个坐标的投影。]
\end{exercise}

\begin{exercise}\label{exercise:ch01-1-12}
证明命题~\cref{proposition:ch01-乘积带边}（光滑流形的乘积连同一个光滑带边流形是一个光滑带边流形）。
\end{exercise}

\end{document}